\chapter{Institutional Discussion, Critical Limitations, and Future Work}
\label{ch:discussion}

\section{Translational Insights for VJTI Administration}
The empirical and algorithmic findings established in this research provide actionable policy guidelines for university administrators, alumni associations, and faculty placement cells:
\begin{enumerate}
    \item \textbf{Abolish Purely Greedy Matching}: The persistent failure of greedy heuristics to protect peripheral students ($Z_2 = 0.547$ under random, $1.023$ under greedy) and their tendency to induce catastrophic mentor overloading ($Z_3 = 8$) demonstrate that legacy first-come-first-served systems undermine mentorship retention.
    \item \textbf{Adopt Pareto-Aware Compromise Solutions}: Institutional administrators should not seek a single scalar assignment. Instead, by presenting the multi-objective Pareto front generated by the Hybrid Memetic solver, leadership can choose a compromise solution based on administrative priorities:
    \begin{itemize}
        \item \textit{Utilitarian Focus}: Prioritizing high industry placement scores ($Z_1 = 72.62, Z_3 = 6$).
        \item \textit{Workload Sustainability}: Restricting mentor load to $Z_3 = 5$, accepting a negligible 2.1\% loss in total quality.
    \end{itemize}
    \item \textbf{Dynamic Threshold Setting}: Rather than imposing rigid, untested administrative rules, institutions should adopt the data-driven calibration framework ($\theta = 0.85 Z_2^{\text{greedy}}$, $\delta = 1.20 \Delta_{\text{greedy}}$) to guarantee operational feasibility.
\end{enumerate}

\section{Critical Methodological Limitations}
In the interest of scientific rigor and transparent academic reporting, we explicitly document the limitations of the current study:
\begin{itemize}
    \item \textbf{Synthetic Data Environment}: Due to student privacy mandates, all profiles and compatibility matrices were synthesized. While the generator incorporates empirical departmental distributions and domain affinities, real human preferences may exhibit unmodeled correlations, idiosyncratic constraints, and subjective preferences.
    \item \textbf{Static Preference Model}: The formulation assumes fixed, deterministic compatibility scores $q_{ij}$. In real-world environments, mentor and student preferences may be dynamic, incomplete, or revealed sequentially through interview interactions.
    \item \textbf{Single-Run Exact Solvers}: Exact MILP Pareto bounds were established exclusively on the small cohort due to exponential branch-and-cut complexity. Medium and large cohorts rely on evolutionary approximations without guaranteed global optimality certificates.
    \item \textbf{Feasibility Rates in Meta-Heuristics}: On highly constrained instances, random mutation operators in evolutionary search occasionally generate infeasible offspring (feasibility rate $\approx 76\%$). While the repair operator and elitist archiving successfully salvage feasible Pareto fronts, developing constraint-preserving genetic operators remains an open theoretical challenge.
\end{itemize}

\section{Roadmap for Future Research}
Building upon the mathematical and computational foundation established in this report, several promising research avenues emerge:
\begin{enumerate}
    \item \textbf{Two-Sided Stable Matching Integration}: Formulating the problem within the game-theoretic framework of the National Resident Matching Program (NRMP) or Gale-Shapley Deferred Acceptance, incorporating two-sided stability constraints alongside multi-objective fairness.
    \item \textbf{Robust and Stochastic Optimization}: Incorporating uncertainty into alumni availability, modeling mentor cancellations, dropout probabilities, and stochastic time commitments using chance-constrained programming.
    \item \textbf{Explainable AI and Fairness Auditing}: Developing interpretable post-hoc explanations for student mentor assignments, allowing students to understand the departmental parity and quality considerations governing their allocations.
    \item \textbf{Pilot Deployment at VJTI}: Implementing an institutional pilot within a single academic department (e.g., Computer Engineering), gathering empirical student feedback, and auditing mentor retention rates over a multi-year longitudinal study.
\end{enumerate}

\section{Conclusion}
This report has formulated, solved, and benchmarked the Multi-Objective Capacitated Bipartite Assignment Problem tailored to the academic ecosystem of Veermata Jijabai Technological Institute (VJTI), Mumbai. By reconciling utilitarian placement quality with Rawlsian egalitarian fairness and mentor workload sustainability under seven operational constraints, we have proven that institutional matching need not sacrifice equity to achieve efficiency. 

The developed \textbf{Hybrid Seeded Memetic NSGA-II} algorithm delivers a statistically validated breakthrough, achieving an 8-fold hypervolume advantage over standard baselines, strictly dominating greedy priority matching, and scaling efficiently to 1,000-student university cohorts. The provided codebase, interactive Streamlit demonstration dashboard, and rigorous statistical reports establish an end-to-end foundation for modern, equitable higher-education operations.

